% !PRINTED 617
\weberpage{617}
%% [page furniture, not body text -- source-confirmed 2026-08-10]
%% \S\,168. Reduktion der Klassengleichung in den Kreisteilungsk\"orpern.

%<<CONT
zu bestimmen. Der erste Faktor ist die Klassenzahl der quadratischen Formen der Ordnung $O'$ [\S\,100, (3)], die wir mit $H'$ bezeichnen. Es ist also
\weq{3}{h = H' (\bar{E} O', \bar{E} A) = H' (O', A)}
(nach \S\,100, 13.).

Es bleibt noch $(O', A)$ zu bestimmen.

Bedeutet $\varrho_1, \varrho_2, \dots, \varrho_\mu$ ein vollst\"andiges Repr\"asentantensystem von $O'$ nach $A$, so m\"ussen die Zahlen
\weq{4}{N(\varrho_1),\ N(\varrho_2)\dots,\ N(\varrho_\mu)}
% [sic] printed N(\varrho_2)\dots, with no comma between N(\varrho_2) and the dots, a punctuation misprint in an otherwise comma-separated list (cf. \varrho_1, \varrho_2, \dots, \varrho_\mu in the preceding sentence)
nach dem Modul $m$ alle verschieden sein, und jede Norm einer Zahl in $O$ ist mit einer dieser Zahlen nach $m$ kongruent. Ist daher $Z$ die Gruppe der rationalen Zahlen, $M$ die Gruppe der Zahlen $z$ aus $Z$, die der Bedingung
\[
z \equiv 1 \pmod{m}
\]
gen\"ugen, $R_m$ die Gruppe der Normenreste nach $m$, so ist (4) ein vollst\"andiges Repr\"asentantensystem von $R_m$ nach $M$ und folglich
\[
\mu = (O', A) = (R_m, M).
\]

Daf\"ur kann man auch setzen (\S\,100, 12.):
\weq{5}{(O', A) = \frac{(Z, M)}{(Z, R_m)}.}

Nun ist
\weq{6}{(Z, M) = \varphi(m) = m \Pi\Bigl(1 - \frac{1}{r}\Bigr),}
wenn $r$ die Primfaktoren von $m$ durchl\"auft und folglich
\weq{7}{h = \frac{H'\varphi(m)}{(Z, R_m)}.}

Den Klassenk\"orper $\mathfrak{K}(A)$ unserer Gruppe erhalten wir, wenn wir dem Klassenk\"orper $\mathfrak{K}(D)$ eine primitive $m$te Einheitswurzel adjungieren:
\weq{8}{\mathfrak{K}(A) = \mathfrak{K}(D, \varrho).}
Denn die Bedingung daf\"ur, da\ss\ ein in einer Primzahl $p$ aufgehendes Primideal $\mathfrak{P}$ des K\"orpers $\mathfrak{K}(D,\varrho)$ vom ersten Grade sei, ist, wenn $(k)$ eine Klasseninvariante der Diskriminante $D$ bedeutet:
\weq{9}{(k)^p \equiv (k), \quad \varrho^p \equiv \varrho \pmod{\mathfrak{P}},}
wenn alle st\"orenden Primzahlen, z.\,B. die in den Diskriminanten der die Zahlen $(k)$ und $(\varrho)$ definierenden Gleichungen, in den Exkludenten $S$ aufgenommen sind.
